\subsection{Código BCD (Binary Coded Decimal)}

Embora o sistema binário seja eficiente para processamento, ele não é naturalmente adequado para representação direta de números decimais na forma utilizada por humanos.

Para contornar essa limitação, utiliza-se o código BCD (Binary Coded Decimal), no qual cada dígito decimal (0 a 9) é representado individualmente por um grupo de 4 bits.

Por exemplo, o número decimal $45_{10}$ em BCD é representado como:

\[
45_{10} = 0100\ 0101_{BCD} =  101101_2
\]

No BCD, apenas combinações de 0000 a 1001 são válidas (0 a 9). Valores de 1010 a 1111 são considerados inválidos.